\scanpage{31}
\mspage{- 213 -}\hfill n\textsuperscript{o} 391

11. \uline{Système de type} $E_7$\unsure{.|:}

(I) Gardons les notations du n\textsuperscript{o}~10, et considérons l'intersection $R'$ de $R$ et de l'hyperplan $V'$ de $V$ orthogonal au vecteur $\varepsilon_7+\varepsilon_8$ ; ce sont donc les vecteurs $x=\displaystyle\sum_{i=1}^{8}\xi_i\varepsilon_i$ de $R$ tels que $\xi_7+\xi_8=0$, autrement dit :

\[
\begin{aligned}
&\pm\varepsilon_i\pm\varepsilon_j,\quad(i\neq j,\ 1\leq i\leq6,\ 1\leq j\leq6),\ \pm(\varepsilon_7-\varepsilon_8)\text{ et}\\
&\pm\frac{1}{2}\bigl(\varepsilon_7-\varepsilon_8+\sum_{i=1}^{6}(-1)^{\unsure{\nu|\partial|\mathcal D}(i)}\varepsilon_i\bigr),\quad\text{où }\sum_{i=1}^{6}\nu(i)\text{ est \uline{impair}.}
\end{aligned}
\]

Il est clair que $R'$ engendre $V'$, donc est un système de racines réduit de rang 7, ayant $n=2+\binom{6}{2}.4+2^6=126$ éléments.

(II) Munissons $V'$ de l'ordre induit par l'ordre lexicographique sur $V$ ; les racines de $R'$ qui sont positives pour cet ordre sont, comme on l'a vu au n\textsuperscript{o}~10, de la forme $\displaystyle\sum_{i=1}^{8}n_i\alpha_i$, où les $n_i$ sont des entiers $\geq0$ ; si l'on exprime qu'une telle racine est orthogonale à $\varepsilon_7+\varepsilon_8$, il vient, (puisque la plus grande racine $\varepsilon_7+\varepsilon_8$ de $R$ est orthogonale à tous les $\alpha_i$ sauf $\alpha_8$), $n_8=0$. Toutes les racines positives de $R'$, (pour l'ordre considéré), sont donc combinaisons linéaires à coefficients positifs des $\alpha_i$ pour $1\leq i\leq7$, ce qui montre que ces racines, (qui sont orthogonales à $\varepsilon_7+\varepsilon_8$, donc appartiennent à $R'$), forment un système de \uline{racines simples} $S'$ de $R'$. Le calcul des $(\alpha_i\mid\alpha_j)$ fait au n\textsuperscript{o}~10 montre alors que le schéma de Dynkin de $R'$ est de type $E_7$, et $R'$ est irréductible.

(III) En vertu du \S~8, n\textsuperscript{o}~3, formule (4), on a $h=n/\ell=18$.

(IV) La plus grande racine pour l'ordre induit par celui de $R^8$ défini au n\textsuperscript{o}~10 est maintenant $\varepsilon_8-\varepsilon_7=2\alpha_1+2\alpha_2+3\alpha_3+4\alpha_4+3\alpha_5+2\alpha_6+\alpha_7$, donc c'est la plus grande racine $\tilde\alpha$, (\S~6, n\textsuperscript{o}~6, cor. 2 de la prop. 13), et on vérifie que

% archive sheet number: 30
